\section{Introduction}\label{sec:intro}

An always-on recogniser listens to every second of audio, speech or not, for as long as the
device or service runs. Its cost is paid continuously: the share of a CPU core it occupies, the
watts that share draws, the memory it holds, and how late its words arrive. Anyone choosing a
recogniser for such a service wants these numbers, and its accuracy, on one axis.

Published comparisons do not give them. Each uses its own machine, thread count, audio and text
normaliser: one vendor benchmark averages four corpora on ten cores of a desktop
CPU~\cite{picovoice2026benchmark}; a runtime's own benchmark reports eight threads on another
desktop CPU and a hundred test utterances~\cite{nvidia2026nemospeechbench}; a model card
reports eight cores of a server CPU on a test set~\cite{moondream2026parakeetredux}; a model
family's documentation reports response time on a laptop~\cite{moonshine2026docs}. None of
these numbers can be divided by another.

% Energy column: generated R7 macros (tables.py), one method for every row.
% The measured comparison (TABLE I): Isosyll live and final vs the public CPU
% comparators on the 730-utterance LibriSpeech DEV panel.
\begin{table*}[t]
\centering
\caption{\textbf{Measured comparison on the 730-utterance LibriSpeech DEV panel} (400 dev-clean,
330 dev-other; one scorer and normaliser). Speed: median of three cold runs on one pinned Zen~2
core ($\times$ real time, model load included). RSS: peak resident memory of the process tree (MB).
Artifact: model and LM files loaded (MB). Core~J: R7 matched campaign: idle-subtracted core J per
audio-second, \RsevenClips-clip subset, median of \RsevenCleanMin--\RsevenCleanMax{} core-clean blocks
(not board or wall power); \ultra{} from the R8 full-panel blocks under the same protocol.
$^{\ddagger}$Shared-host timing (interim, host H0); WER is from the full 730-utterance panel.
$^{\S}$Timed on a 73-utterance subset of the panel; WER is from the full 730-utterance panel.
Parakeet TDT v3 is offline (whole utterance). Cell colour: green is better, red worse within a column. DEV-only.}
\label{tab:measured}
\footnotesize
\setlength{\tabcolsep}{3.5pt}
\begin{tabular}{@{}llrrrrrrr@{}}
\toprule
 & & Speed & \multicolumn{3}{c}{Panel WER (\%)} & RSS & Artifact & Core \\
\cmidrule(lr){4-6}
System & Runtime & ($\times$ RT) & clean & other & pooled & (MB) & (MB) & (J) \\
\midrule
\rawinput{data/table_measured.tex}
\bottomrule
\end{tabular}
\end{table*}
%
\para{What we did.} We ran the recognisers of those comparisons that we could obtain and license on
our own hardware, under one protocol (\cref{sec:protocol}): the same 730 LibriSpeech DEV utterances (\ColdAudioMinutes{}
minutes), one cold process per system with model load included, one pinned physical core, three
repetitions, and one scorer and text normaliser for every transcript. We keep two classes apart. A \emph{live streaming} system emits words while the
audio arrives, so it has a word lag. A \emph{whole-utterance} system sees the whole clip before
it emits anything, so it has throughput but no live words. \Cref{fig:money} plots speed against
dev-clean WER for every system measured so far, and \cref{tab:measured} adds dev-other and pooled
WER, memory, word lag, watts and an eight-core budget. Numbers that others published for the same
systems are collected in \cref{tab:reported}; they never share an axis with a measured number.

\para{The cost of always listening.} \Cref{tab:alwayson} gives, for every system measured on the
quiet host, the share of one core, the watts and the kilowatt-hours per year one always-on stream
costs, its memory and its live word lag. \ColdWattsSentence{}

\para{Where \ultra wins and where it loses.} \ColdHeadSentence{} \ColdSpeedSentence{}
\ColdAccuracySentence{} \ColdMemorySentence{} \ultra is trained on LibriSpeech alone; training
continues, and the engine does not change with the weights.

\para{Our approach.} \ultra co-designs a small streaming transducer and a CPU engine around one
invariant: \emph{integer exactness}. Every quantised GEMM accumulates exactly in int32; every
activation row is quantised from its own range; and the log-mel front end is computed in
portable IEEE arithmetic. The engine's tokens are therefore a pure function of the input audio
and the weights, independent of batch composition, kernel tile and instruction set. Exactness
lets us change kernels freely (any change must reproduce the continuity fixtures bit for bit)
and makes every speed measurement a measurement of the same computation.

\para{Units.} CPU cost is CPU-seconds per audio-second (\cpus), the measured work of the whole
process tree. Speed is its reciprocal, audio-seconds per CPU-second, written as a multiple of
real time; on one pinned core it is also the number of live streams that core could feed.
Density is a paced pass under a lateness rule. We never convert one of these into another
across harnesses.

\para{Two cost boundaries.} \emph{Cold} cost times a whole recogniser process, model load
included; any recogniser that runs as a process can be measured this way, so the comparison of
\cref{fig:money} uses it. \emph{Warm} cost times a running service after load, inside the one
harness that also measures paced density (\cref{sec:density}). The two are never mixed.

\para{Contributions.}
\begin{enumerate}
\item \textbf{A same-hardware comparison of published CPU recognisers}: one protocol for speed,
  WER, memory, word lag and energy, with live streaming and whole-utterance systems kept apart
  and published numbers kept in a separate group (\cref{fig:money,tab:measured,tab:reported}).
\item \textbf{A batch-invariant, cross-ISA bit-exact INT8 engine.} Per-row centred activations
  with per-channel weights cost at most \GateMaxDelta\pp WER against FP32, are bit-identical at
  batch sizes 1 to 32, and give \XisaAgree/\XisaTotal{} identical PCM-to-token decodes on AVX2 and
  NEON (\cref{sec:engine}).
\item \textbf{A multirate streaming encoder} running ten of sixteen blocks at \MidRateHz\,Hz inside
  160\,ms chunks, and a paced serving-density measurement: \DenMRCore{} live streams per core at
  a worst-worker p99 lateness of 800\,ms against \DenRzeroCore{} for Zipformer INT8 in the same
  harness, each beside its own checkpoint's WER (\cref{sec:model,tab:density}).
\item \textbf{A pre-registered evaluation protocol} for CPU ASR: a single harness, a fixed density
  ladder that stops at the first failure, counterbalanced timing blocks with Holm correction, a
  speaker-balanced selection panel, and a bootstrap INT8 gate (\cref{sec:protocol}).
\item \textbf{Negative results}: ternary weights that fit in L3 are \emph{not} faster at batch on a
  compute-bound CPU, and an unsigned 7-bit GEMM tile that removes three instructions is slower on
  Zen~2 (\cref{sec:negative}).
\end{enumerate}

\para{What we do not claim.} \ultra is less accurate than the best systems it outruns, and every
speed in this paper is printed next to the WER of the same checkpoint, harness and run. Our
accuracy comes from DEV data only; the LibriSpeech test sets stay sealed. All results are
exploratory; qualified numbers will replace them through the same macros once the protocol is
frozen.
